\chapter{Literature Review}
\label{chap:lr}

This chapter reviews backdoor attacks, defenses against them, perceptual hash
functions and MPC for machine learning, and then states the research gap.

\section{Backdoor Attacks}
\label{sec:lr-attacks}

Among the poisoning attacks studied since the work of Biggio
\emph{et al.}~\cite{biggio2012poisoning}, backdoors are the most insidious,
because the poisoned model keeps its accuracy on clean
data~\cite{goldblum2023dataset}. Li
\emph{et al.}~\cite{li2024backdoorsurvey} classify backdoor attacks along several
axes, four of which decide whether a data-level defense can work: the adversary
may only poison the data or also control training; the poisoned samples may be
relabelled (dirty-label) or keep their labels (clean-label); the trigger may be
visible or invisible; and it may be the same for every sample
(sample-agnostic) or differ between samples (sample-specific). A filter applied
to the training data is meaningful only against a poison-only adversary.

BadNets~\cite{gu2019badnets} stamped a small fixed patch onto the poisoned
images, and the Blended attack~\cite{chen2017blended} mixed a full-image pattern
into them with a weight $\alpha$ that can be made small enough for the trigger to
become invisible. Clean-label attacks keep the correct labels:
Turner \emph{et al.}~\cite{turner2019labelconsistent} combined a patch with
adversarial perturbations, Saha \emph{et al.}~\cite{saha2020hidden} and Souri
\emph{et al.}~\cite{souri2022sleeper} hid the trigger completely at training
time, and Refool~\cite{liu2020refool} used natural reflections. Sample-specific
invisible triggers include the warping field of WaNet~\cite{nguyen2021wanet} and
the encoder-generated residual of ISSBA~\cite{li2021issba}, whereas
IAD~\cite{nguyen2020iad}, LIRA~\cite{doan2021lira} and blind
backdoors~\cite{bagdasaryan2021blind} require control over training and lie
outside the scope of any data filter. BATT~\cite{xu2023batt} uses a spatial
transformation as the trigger. Finally, Qi
\emph{et al.}~\cite{qi2023revisiting} designed adaptive attacks with several
low-opacity triggers and correctly labelled cover samples, which removed the
separation between poisoned and clean samples in the latent space of the model.

Attacks are evaluated by the benign accuracy (BA) on clean test data and by the
attack success rate (ASR), the share of triggered test inputs classified as the
target class~\cite{li2024backdoorsurvey}. Purification defenses are additionally
evaluated by the area under the ROC curve (AUC)~\cite{fawcett2006roc} of their
suspicion score and by the shares of poisoned and clean samples they remove.

\section{Backdoor Defenses}
\label{sec:lr-defenses}

Existing defenses act at different stages of the pipeline~\cite{li2024backdoorsurvey,li2023backdoorbox}.
Input pre-processing~\cite{liu2017neuraltrojans,li2021shrinkpad} and input-level
detection~\cite{gao2019strip,guo2023scaleup,hou2024ibdpsc} operate on a deployed
model at inference time. Model repair methods such as
Fine-Pruning~\cite{liu2018fineprune}, Neural Attention
Distillation~\cite{li2021nad} and Neural Cleanse~\cite{wang2019neuralcleanse}
require a trained model and a reserved set of clean data. Poison suppression
methods~\cite{li2021abl,huang2022dbd} modify the training procedure, and
certified defenses such as Deep Partition Aggregation~\cite{levine2021dpa} train
many models on partitions of the data.

Dataset purification removes poisoned samples before training and needs no
clean data. Spectral Signatures~\cite{tran2018spectral} scored the
samples of each class by their projection onto the top singular vector of the
penultimate-layer features, and Activation Clustering~\cite{chen2019activation}
split the same features into two clusters with $k$-means.
SPECTRE~\cite{hayase2021spectre} and SCAn~\cite{tang2021scan} replaced these
statistics with robust covariance estimation and an EM-based decomposition.
After the adaptive attacks of Qi \emph{et al.}~\cite{qi2023revisiting},
ASSET~\cite{pan2023asset}, confusion training~\cite{qi2023proactive} and
FLARE~\cite{hou2025flare} began to shape the training process or to use all
hidden layers. All of these methods, however, compute their statistic in the
feature space of a network trained on the poisoned data, using SVD, clustering
or EM. Zeng \emph{et al.}~\cite{zeng2021frequency} detect triggers by their
high-frequency artefacts without the victim model, but their detector is itself
a neural network applied to inference-time inputs.

\section{Perceptual Hash Functions}
\label{sec:lr-phash}

A perceptual hash maps an image to a short bit string so that perceptually
similar images are close in Hamming distance, whereas unrelated images are
not~\cite{venkatesan2000robust,monga2006perceptual,zauner2010phash,farid2021overview}.
Classical constructions threshold a downscaled image or its gradients (average
and difference hashes), its low-frequency DCT~\cite{ahmed1974dct} or wavelet
coefficients~\cite{mallat1989wavelet} (pHash, the wavelet hash and
PDQ~\cite{meta2019pdq}) or block statistics~\cite{yang2006blockmean}. Apart
from the colour hash and the proprietary or learned PhotoDNA and
NeuralHash~\cite{apple2021neuralhash}, all of them apply a fixed linear transform
and compare each coefficient with a threshold.

Perceptual hashes are not cryptographic, and they have been attacked
repeatedly~\cite{hao2021notwhat,struppek2022neuralhash,jain2022adversarial,prokos2023squint}:
small, often imperceptible changes suffice to move an image far from its hash or
onto the hash of another image. For this thesis, the lesson is that invariance
cuts both ways: a hash that ignores small changes is also blind to an invisible
trigger. In machine learning pipelines, hashing has been used for
near-duplicate removal~\cite{barz2020cifair}, for integrity checks of web-scale
datasets~\cite{carlini2024webscale} and, under cryptographic protection, for
matching images against known harmful content~\cite{kulshrestha2021identifying}.
To the best of the author's knowledge, it has never been used to filter poisoned
training samples.

\section{Secure Multi-Party Computation for Machine Learning}
\label{sec:lr-mpc}

MPC allows parties to compute a joint function of their private inputs without
revealing them~\cite{yao1986generate,goldreich1987play,benor1988completeness}.
Privacy-preserving machine learning relies mainly on honest-majority
three-party replicated secret sharing~\cite{araki2016highthroughput,mohassel2018aby3,wagh2021falcon}
and on dishonest-majority protocols of the SPDZ
family~\cite{damgard2012spdz,cramer2018spdz2k}, implemented in frameworks such as
SecureML~\cite{mohassel2017secureml}, CrypTen~\cite{knott2021crypten} and
MP-SPDZ~\cite{keller2020mpspdz}. These protocols share a cost structure. Linear
operations with public coefficients, including a DCT or a wavelet transform,
are local and free; a multiplication of two secrets costs one round with
precomputed triples~\cite{beaver1991efficient}; a comparison requires a bit
decomposition with $O(k\log k)$ AND gates and $O(\log k)$ rounds for a $k$-bit
ring~\cite{catrina2010primitives,escudero2020edabits}; and sorting and iterative
algorithms such as $k$-means multiply this cost by the number of comparisons and
iterations~\cite{batcher1968sorting,mohassel2020kmeans}. A defense that first
trains a deep network and then runs SVD or clustering is therefore impractical
under MPC.

Poisoning in privacy-preserving learning has been studied mainly for federated
learning, where a single participant can implant a
backdoor~\cite{bagdasaryan2020backdoorfl} and secure
aggregation~\cite{bonawitz2017secagg} hides individual updates. Robust secure
aggregation schemes~\cite{so2021brea,nguyen2022flame,lycklama2023rofl,rathee2023elsa}
filter or bound model updates with cheap statistics such as norms, and
SafeNet~\cite{chaudhari2023safenet} limits the influence of a minority of
poisoned data owners in MPC training with an ensemble of locally trained
models. None of these defenses filters individual samples.

\section{Summary and Research Gap}
\label{sec:lr-summary}

Table~\ref{tab:lr-gap} compares the reviewed approaches with the requirements of
this thesis. Dataset purification defenses depend on a model trained on the
poisoned data, defenses for privacy-preserving learning do not work at the level
of samples, and perceptual hashes fit the MPC cost model but have never been
evaluated against poisoning. The thesis therefore builds a defense on
per-sample hash bits and on statistics that need no trained model, adds hashes
that keep the high-frequency component where triggers
reside~\cite{zeng2021frequency}, evaluates it with the established metrics
against attacks from BackdoorBox~\cite{li2023backdoorbox}, including adaptive
ones~\cite{qi2023revisiting}, and estimates its cost under three-party and
two-party MPC protocols~\cite{mohassel2018aby3,cramer2018spdz2k}.

\begin{table}[htbp]
\centering
\caption{Comparison of Existing Approaches}
\label{tab:lr-gap}
\setstretch{1.0}
\footnotesize
\begin{tabular}{>{\raggedright\arraybackslash}p{0.30\textwidth} >{\raggedright\arraybackslash}p{0.10\textwidth} >{\raggedright\arraybackslash}p{0.11\textwidth} >{\raggedright\arraybackslash}p{0.11\textwidth} >{\raggedright\arraybackslash}p{0.18\textwidth}}
\toprule
\textbf{Approach} & \textbf{Sample level} & \textbf{No trained model} & \textbf{Cheap under MPC} & \textbf{Invisible triggers} \\
\midrule
Model-based purification~\cite{tran2018spectral,chen2019activation,hayase2021spectre,tang2021scan,hou2025flare} & yes & no & no & yes, weakened by adaptive attacks \\
Input-level detection~\cite{gao2019strip,guo2023scaleup,hou2024ibdpsc} & per query & no & no & yes \\
Robust secure aggregation~\cite{so2021brea,nguyen2022flame,lycklama2023rofl,rathee2023elsa} & no & no & partly & not applicable \\
MPC ensembles~\cite{chaudhari2023safenet} & no & no & yes & not applicable \\
Perceptual hashing & yes & yes & yes & unknown \\
\bottomrule
\end{tabular}
\end{table}
